% ThemisDB: Hardware-Normalized Performance Evaluation of a
% Hybrid Multi-Model Database System
%
% IEEE Conference Format (IEEEtran class)
% Submit with: pdflatex + bibtex + pdflatex + pdflatex
%
\documentclass[conference]{IEEEtran}

% ---- packages ----------------------------------------------------------------
\usepackage[T1]{fontenc}
\usepackage[utf8]{inputenc}
\usepackage{amsmath}
\usepackage{amssymb}
\usepackage{graphicx}
\usepackage{booktabs}
\usepackage{multirow}
\usepackage{tabularx}
\usepackage{array}
\usepackage{url}
\usepackage{hyperref}
\usepackage{xcolor}
\usepackage{cite}
\usepackage{balance}

% ---- column types for tables -------------------------------------------------
\newcolumntype{L}[1]{>{\raggedright\arraybackslash}p{#1}}
\newcolumntype{R}[1]{>{\raggedleft\arraybackslash}p{#1}}
\newcolumntype{C}[1]{>{\centering\arraybackslash}p{#1}}

% ---- macros ------------------------------------------------------------------
\newcommand{\themis}{\textsc{ThemisDB}}
\newcommand{\ecclass}[1]{$E_{\text{#1}}$}
\newcommand{\nfact}[1]{$N_{\text{#1}}$}

% ==============================================================================
\begin{document}

\title{ThemisDB: Hardware-Normalized Performance Evaluation of a
       Hybrid Multi-Model Database System}

\author{%
  \IEEEauthorblockN{ThemisDB Development Team}
  \IEEEauthorblockA{\textit{makr-code/ThemisDB} \\
    \url{https://github.com/makr-code/ThemisDB}}
}

\maketitle

% ==============================================================================
\begin{abstract}
We present a longitudinal performance evaluation of \themis{}, a
C++20 hybrid multi-model database that integrates graph, vector,
relational, time-series, and document data models within a single
ACID-compliant system governed by an AQL query engine and
multi-version concurrency control (MVCC).
The study spans versions v1.3.0 through v1.8.2 and covers
1\,078+ benchmark cases distributed across 33 software modules.
Because absolute throughput numbers are platform-dependent and
therefore non-transferable across hardware configurations, we
introduce a \emph{hardware-normalized efficiency model} that
decomposes the expected throughput of each workload class into a
weighted linear combination of eight normalized hardware factors:
CPU integer throughput, memory bandwidth (STREAM-triad), sequential
and random storage I/O, GPU VRAM capacity, and PCIe host-to-device /
device-to-host transfer rates plus dispatch latency.
We validate the model on a 20-core x86-64 development platform and
report hardware-neutral scores for the six benchmark classes defined
in the CHIMERA test suite, which is compliant with IEEE~Std~2807-2022
and ISO/IEC~14756:2015.
Key findings confirm that PCIe transfer bandwidth and memory bandwidth
are the dominant bottlenecks for GPU-heavy workloads, while random
I/O efficiency already meets target levels on the reference platform.
All v0 efficiency values exceed 1.0, indicating conservative
baselines that will be recalibrated as multi-host data accumulates.
\end{abstract}

\begin{IEEEkeywords}
multi-model database, performance evaluation, hardware normalization,
benchmark, vector search, LLM inference, time-series, graph database,
AQL, MVCC
\end{IEEEkeywords}

% ==============================================================================
\section{Introduction}
\label{sec:intro}

Modern application stacks increasingly require databases that
simultaneously support relational queries, graph traversals, vector
similarity search, and time-series analytics~\cite{Lu2018,Stonebraker2010}.
Separate purpose-built systems impose heavy operational overhead and
cross-store consistency burdens.
Polystore approaches~\cite{Atzeni2016} address data variety at the
abstraction level but typically do not provide native transactional
guarantees across models.

\themis{} is an open-source, C++20 hybrid multi-model database system
that provides graph, vector, relational, time-series, and document
models inside a single process with full ACID support.
Its AQL query engine, adapted from ArangoDB's query language,
compiles multi-paradigm queries — including LLM inference, RAG
retrieval, and geospatial joins — into a single execution plan.
The storage layer wraps RocksDB~\cite{Dong2021} for write-ahead
logging, compaction, and point-in-time recovery, while the
replication layer employs Raft consensus~\cite{Ongaro2014}.

\textbf{Contributions of this paper:}
\begin{enumerate}
  \item \emph{Longitudinal benchmark study} across five minor release
        series (v1.3.0--v1.8.2) covering key performance indicators
        (KPIs) for query, vector, graph, and time-series workloads
        using the Google Benchmark framework~\cite{GoogleBenchmark}.
  \item \emph{Hardware-normalized efficiency model} that decomposes
        expected throughput into eight weighted hardware factors and
        produces a platform-independent \emph{score\_hw\_neutral}
        for each benchmark run.
  \item \emph{Governance framework} with mandatory root-cause analysis
        when score\_hw\_neutral $< 0.90$ across two consecutive runs,
        and automatic recalibration proposals when the score exceeds
        $1.10$ across three or more distinct hardware platforms.
\end{enumerate}

% ==============================================================================
\section{System Architecture Overview}
\label{sec:arch}

\themis{} is organized in seven horizontal layers, each encapsulating
a distinct concern.
Table~\ref{tab:arch} summarizes the layers, selected modules, and
key technologies.
Across these layers the system comprises 46 software modules, 153
documentation files, and a test suite of 1\,078+ benchmark cases plus
over 1\,200 unit/integration tests (v1.8.2, partial run).

\begin{table}[t]
  \caption{ThemisDB Architecture Layers (46 modules)}
  \label{tab:arch}
  \centering
  \begin{tabular}{L{1.5cm}L{2.6cm}L{3.5cm}}
    \toprule
    \textbf{Layer} & \textbf{Modules} & \textbf{Key Technologies} \\
    \midrule
    Query \& Index &
      AQL Parser, Optimizer, JIT, CTE &
      Multi-paradigm query rewrite, B-Tree, R-Tree, HNSW
      (CPU \& GPU), secondary index \\
    \addlinespace
    Storage &
      RocksDB WAL, Compaction, PITR, Blob Backends &
      Seven object backends (S3/Azure/GCS/MinIO/local/
      memory/hybrid), snapshot isolation \\
    \addlinespace
    Analytics &
      OLAP, CEP, IVM, Streaming Aggregation &
      CUBE/ROLLUP, SIMD vectorization (4.5$\times$--6.9$\times$
      speedup), process mining \\
    \addlinespace
    Time-Series &
      TS Storage, Gorilla Encoder, Retention &
      Adaptive flush, Gorilla delta-of-delta
      encoding~\cite{Pelkonen2015}, tiered retention \\
    \addlinespace
    Graph &
      Graph Index, Traversal Engine &
      Adjacency index, BFS/DFS/Dijkstra/A*/Bidirectional,
      12 constraint types \\
    \addlinespace
    Vector / AI &
      ANN Search, Embedding Cache, LLM, RAG, Training &
      HNSW~\cite{Malkov2018} GPU search, llama.cpp
      inference~\cite{llamacpp2023}, LoRA adapters~\cite{Hu2022},
      RAG~\cite{Lewis2020} pipeline \\
    \addlinespace
    Distributed &
      Replication, Sharding, 2PC, SAGA &
      Raft~\cite{Ongaro2014}, semi-sync/async replication,
      2PC~\cite{Gray1978}, SAGA~\cite{Garcia-Molina1987} \\
    \bottomrule
  \end{tabular}
\end{table}

The security layer provides AES-256-GCM field-level encryption, TLS
1.3 termination, RBAC, and SOC~2 / NIST / GDPR compliance hooks.
Seven network protocols are supported: HTTP/1.1, HTTP/2, HTTP/3,
WebSocket, MQTT, PostgreSQL Wire, and gRPC.

\subsection{Module-Level SLO Inventory}
\label{ssec:slo-inventory}

Each module carries explicit service-level objectives (SLOs)
maintained in per-module \texttt{FUTURE\_ENHANCEMENTS.md} files.
Table~\ref{tab:slo-summary} lists representative SLOs across
14 selected modules together with the measurement status in v1.8.2.
A total of 28 modules define at least one SLO; of these, 12 have
at least one directly measured value, 10 have only proxy coverage,
and 6 lack any current measurement (no benchmark runner activated).

\begin{table}[t]
  \caption{Representative Module SLOs and v1.8.2 Status}
  \label{tab:slo-summary}
  \centering
  \renewcommand{\arraystretch}{1.1}
  \begin{tabular}{L{1.2cm}L{1.3cm}R{1.3cm}L{1.0cm}L{1.5cm}}
    \toprule
    \textbf{Mod.} & \textbf{SLO ID} & \textbf{Target} &
    \textbf{v1.8.2} & \textbf{Status} \\
    \midrule
    Query   & Q-1 & 750\,M\,items/s & 796\,M & \checkmark \\
    Index   & I-1 & 280\,k\,ins/s  & 549\,k & \checkmark \\
    Index   & I-2 & 180\,k\,ins/s  & 255\,k & \checkmark \\
    Cache   & C-1 & 5\,M\,ops/s    & 5.9\,M & \checkmark \\
    Storage & S-1 & 100\,k\,ops/s  & $\sim$1.3\,k & $\times$ \\
    Graph   & G-1 & 500\,k\,e/s    & 1.18\,M & \checkmark \\
    TS      & TS-1 & 500\,k\,pts/s  & 61\,M* & \checkmark \\
    TS      & TS-2 & 2\,GB/s/core   & n/a & -- \\
    Tx      & TX-3 & 6\,k\,2PC/s    & 6.4\,k & \checkmark \\
    Replica & R-1  & P99 $\leq$50ms  & n/a & -- \\
    LLM     & L-1  & TTFT$\leq$200ms & n/a & -- \\
    Search  & SE-1 & P99$\leq$20ms   & n/a & -- \\
    Auth    & AUT-4 & $\leq$1\,$\mu$s (Bloom) & n/a & -- \\
    CDC     & CDC-2 & P99$\leq$20ms  & n/a & -- \\
    \bottomrule
  \end{tabular}
  \vspace{0.5ex}
  {\footnotesize * In-memory append; sustained NVMe write not yet measured separately.
    $\times$ = SLO not met; -- = no current measurement.}
\end{table}

% ==============================================================================
\section{Benchmark Methodology}
\label{sec:method}

\subsection{Benchmark Infrastructure}
\label{ssec:infra}

All micro-benchmarks are implemented with the Google Benchmark
framework~\cite{GoogleBenchmark}.
Each fixture applies at least two warmup iterations followed by a
minimum of five measurement samples; the median is used as the
standard reported value.
Reporting units follow the workload type: operations per second
(ops/s) for query and graph workloads, insertions per second (k/s or
M/s) for index and vector workloads, and points per second (M~pts/s)
for time-series workloads.

\textbf{Platform specifications:}
\begin{itemize}
  \item \emph{v1.3.0--v1.3.3} (baseline runs 20251223,
        20251223\_085556): Intel i9-10900K, 10C/20T \@ 3.70\,GHz,
        31\,GB RAM, WSL2 Linux, MSVC Release x64, AVX2.
  \item \emph{v1.3.4 and v1.8.2} (run 20251229\_184507 and current):
        Windows x64, 20 cores \@ 3.696\,GHz, 20\,MB L3-cache,
        L1=32\,KB, L2=256\,KB, MSVC Release x64, AVX2.
\end{itemize}

\subsection{Hardware Baseline Capture}
\label{ssec:hw-baseline}

Every benchmark run is preceded by the
\texttt{HardwareBaseline.CaptureAndPersist} Google Test fixture, which
writes a structured JSON snapshot to
\texttt{logs/hardware\_baseline/}.
The snapshot schema is summarized in Table~\ref{tab:hw-schema}.

\begin{table}[t]
  \caption{Hardware Baseline Snapshot Schema}
  \label{tab:hw-schema}
  \centering
  \begin{tabular}{L{2.0cm}L{5.5cm}}
    \toprule
    \textbf{Field group} & \textbf{Contents} \\
    \midrule
    Identification &
      \texttt{schema\_version}, \texttt{run\_id},
      \texttt{host\_fingerprint} \\
    Build context &
      compiler, optimization flags, AVX tier \\
    OS context &
      OS name, kernel version, NUMA topology \\
    CPU metrics &
      Integer ops/s (STREAM-like integer kernel) \\
    Memory metrics &
      STREAM triad bandwidth (GB/s) \\
    Storage metrics &
      Sequential read (MB/s), random read IOPS
      (\texttt{fio}-equivalent) \\
    GPU inventory &
      VRAM (GB), H2D transfer (GB/s), D2H transfer (GB/s),
      dispatch latency ($\mu$s) \\
    Runtime context &
      process affinity, huge-page status \\
    \bottomrule
  \end{tabular}
\end{table}

Missing hardware metrics are recorded as \texttt{unavailable} (not
as 0.0) to prevent silent misattribution in subsequent computations.
When a factor is unavailable, it is imputed to 0.5 for the relevant
efficiency formula and flagged as uncertain.

\subsection{CHIMERA Benchmark Suite}
\label{ssec:chimera}

\themis{} integrates the CHIMERA vendor-neutral benchmark adapter
(\texttt{src/chimera}), compliant with IEEE~Std~2807-2022~\cite{IEEE2807}
and ISO/IEC~14756:2015~\cite{ISO14756}.
The adapter provides unified drivers for YCSB~\cite{Cooper2010},
TPC-C~\cite{TPCC}, TPC-H~\cite{TPCH}, ANN-Benchmarks~\cite{AnnBenchmarks},
LDBC-SNB~\cite{Angles2018,LDBCSNB}, vLLM inference, and RAGBench.
CHIMERA normalizes results across participating database systems to
enable apples-to-apples comparisons independent of vendor-specific
reporting conventions.

% ==============================================================================
\section{Hardware-Normalized Efficiency Model}
\label{sec:hwmodel}

\subsection{Problem Statement}
\label{ssec:prob}

Absolute throughput targets depend on the platform on which the
measurement was taken.
A database achieving 796\,M\,ops/s on a 20-core workstation with
31\,GB RAM cannot be directly compared to the same target on a
dual-socket server with 256\,GB RAM and NVMe RAID storage.
A hardware-agnostic evaluation framework requires that expected
throughput be expressed as a function of the actual hardware
capability of the test host.

\subsection{Normalized Hardware Factors}
\label{ssec:norm-factors}

We define eight normalized hardware factors
$N_i \in [0, 1]$ by clamping measured values against reference
ceilings chosen to represent high-end commodity hardware:

\begin{align}
  N_{\text{cpu}}      &= \min\!\bigl(\text{cpu\_ops} / 8{\times}10^{7},\ 1\bigr) \label{eq:ncpu}\\
  N_{\text{mem}}      &= \min\!\bigl(\text{triad\_GB\_s} / 35,\ 1\bigr) \label{eq:nmem}\\
  N_{\text{seq}}      &= \min\!\bigl(\text{disk\_MB\_s} / 1200,\ 1\bigr) \label{eq:nseq}\\
  N_{\text{rand}}     &= \min\!\bigl(\text{disk\_IOPS} / 50000,\ 1\bigr) \label{eq:nrand}\\
  N_{\text{vram}}     &= \min\!\bigl(\text{vram\_GB} / 16,\ 1\bigr) \label{eq:nvram}\\
  N_{\text{h2d}}      &= \min\!\bigl(\text{h2d\_GB\_s} / 8,\ 1\bigr) \label{eq:nh2d}\\
  N_{\text{d2h}}      &= \min\!\bigl(\text{d2h\_GB\_s} / 8,\ 1\bigr) \label{eq:nd2h}\\
  N_{\text{dispatch}} &= \min\!\bigl(200 / \text{dispatch\_us},\ 1\bigr) \label{eq:ndisp}
\end{align}

\subsection{Expected Capacity Models}
\label{ssec:expected}

Six workload classes are defined, each with a weighted linear
expected-capacity function $E_c$ (Table~\ref{tab:emodels}).
Weights reflect the dominant resource bottleneck for each class and
were chosen from first-principles analysis of the system's hot paths;
statistical calibration is planned once 30+ paired runs across 3+
hardware classes are available (Section~\ref{ssec:calibration}).

\begin{table}[t]
  \caption{Expected-Capacity Models per Workload Class (v0)}
  \label{tab:emodels}
  \renewcommand{\arraystretch}{1.2}
  \centering
  \begin{tabular}{L{1.9cm}L{5.7cm}}
    \toprule
    \textbf{Class} & \textbf{Expected capacity $E_c$} \\
    \midrule
    OLTP &
      $0.45\,N_{\text{cpu}} + 0.20\,N_{\text{mem}}
        + 0.25\,N_{\text{rand}} + 0.10\,N_{\text{seq}}$ \\
    OLAP &
      $0.40\,N_{\text{cpu}} + 0.45\,N_{\text{mem}}
        + 0.15\,N_{\text{seq}}$ \\
    Storage WAL &
      $0.20\,N_{\text{cpu}} + 0.20\,N_{\text{mem}}
        + 0.35\,N_{\text{seq}} + 0.25\,N_{\text{rand}}$ \\
    Vector GPU &
      $0.20\,N_{\text{cpu}} + 0.15\,N_{\text{mem}}
        + 0.30\,N_{\text{vram}} + 0.20\,N_{\text{h2d}}
        + 0.10\,N_{\text{d2h}} + 0.05\,N_{\text{dispatch}}$ \\
    LLM GPU &
      $0.15\,N_{\text{cpu}} + 0.15\,N_{\text{mem}}
        + 0.35\,N_{\text{vram}} + 0.20\,N_{\text{h2d}}
        + 0.05\,N_{\text{d2h}} + 0.10\,N_{\text{dispatch}}$ \\
    Mixed CPU+GPU &
      $0.25\,N_{\text{cpu}} + 0.20\,N_{\text{mem}}
        + 0.20\,N_{\text{vram}} + 0.15\,N_{\text{h2d}}
        + 0.10\,N_{\text{d2h}} + 0.10\,N_{\text{dispatch}}$ \\
    \bottomrule
  \end{tabular}
\end{table}

\subsection{Efficiency Computation}
\label{ssec:eff}

For a specific benchmark $b$ assigned to workload class $c$:

\begin{align}
  \text{expected}_{b}  &= \text{baseline\_ref}_{b} \times E_c \\
  \text{efficiency}_{b} &= \frac{\text{measured}_{b}}{\text{expected}_{b}} \\
  \text{score\_hw\_neutral}_{b} &=
        \frac{\text{measured}_{b}}{\text{target\_product}_{b} \times E_c}
\end{align}

\noindent
where \(\text{baseline\_ref}_{b}\) is the reference throughput value
fixed at v0 initialization and \(\text{target\_product}_{b}\) is the
product-roadmap absolute target.
\(\text{score\_hw\_neutral}_{b}\) is the primary release-gate metric.

\subsection{Calibration Rules (v1)}
\label{ssec:calibration}

Version~1 calibration is binding once 30 paired runs across at least
three distinct hardware classes are collected:

\begin{enumerate}
  \item \emph{Class calibration factor:}
        $K_c = \operatorname{median}(\{\text{efficiency}_{b}\}_{b \in c})$
        over the reference window.
  \item \emph{Calibrated efficiency:}
        $\text{efficiency\_cal}_{b} = \text{efficiency}_{b} / K_c$.
  \item \emph{Factor significance:} a hardware factor $N_i$ is deemed
        significant for class $c$ only if
        $|\rho_{\text{Spearman}}(N_i,\,\text{efficiency\_cal})| \geq 0.35$
        and $p \leq 0.05$.
  \item \emph{Stability rule:} the coefficient of variation
        $\text{CV} \leq 0.20$ within the reference window.
\end{enumerate}

\noindent
\textbf{Banding table:}
\begin{itemize}
  \item $\text{score} < 0.85$: \emph{critical}
  \item $0.85 \leq \text{score} < 0.95$: \emph{notable}
  \item $0.95 \leq \text{score} \leq 1.05$: \emph{normal}
  \item $\text{score} > 1.05$: \emph{above target}
\end{itemize}

% ==============================================================================
\section{Experimental Results}
\label{sec:results}

\subsection{Longitudinal Throughput Evolution}
\label{ssec:longitudinal}

Table~\ref{tab:longitudinal} tracks six core KPIs from v1.3.0 through
v1.8.2 together with the v1.8.0 product target.
The benchmark count grew from 450 to 1\,078~(+140\%), reflecting both
new module coverage and test-suite consolidation.

\begin{table*}[t]
  \caption{Longitudinal Throughput Evolution (v1.3.0 -- v1.8.2)}
  \label{tab:longitudinal}
  \centering
  \begin{tabular}{L{2.8cm}R{1.5cm}R{1.5cm}R{1.5cm}R{1.5cm}R{1.6cm}R{1.6cm}R{1.6cm}}
    \toprule
    \textbf{Metric} &
    \textbf{v1.3.0} & \textbf{v1.3.1} & \textbf{v1.3.2} &
    \textbf{v1.3.3} & \textbf{v1.3.4} & \textbf{v1.8.2} &
    \textbf{Target} \\
    \midrule
    Query Engine (M\,ops/s)
      & 700 & 750 & 800 & 800 & \textbf{814.5} & \textbf{796.4} & 900 \\
    Vector Insert (k/s)
      & 280 & 300 & 330 & 340 & \textbf{351.4} & \textbf{548.7} & 600 \\
    Secondary Index (k/s)
      & 180 & 190 & 210 & 215 & \textbf{217.2} & \textbf{254.9} & 1\,000 \\
    Graph Edge Ops (k/s)
      & --- & --- & --- & --- & \textbf{628.7} & \textbf{1\,177} & 1\,000 \\
    Timeseries Insert (M\,pts/s)
      & --- & --- & --- & --- & \textbf{49.0} & \textbf{61.0} & 60 \\
    Benchmark Tests
      & 450 & 480 & 520 & 780 & \textbf{1\,078} & 5\,(partial) & 1\,200 \\
    \bottomrule
  \end{tabular}
\end{table*}

Graph edge operations surpassed the 1\,M\,edges/s target already in
v1.8.2 ($+87\%$ over v1.3.4).
Time-series insert throughput exceeded the 60\,M\,pts/s target by
1.7\%.
Query engine throughput reached 814.5\,M\,ops/s in v1.3.4 but
regressed to 796.4\,M\,ops/s in v1.8.2 due to increased
governance-framework overhead in the hot evaluation path.
Secondary index insert performance remains the largest gap at
254.9\,k/s vs.\ the 1\,M/s target, primarily attributable to
RocksDB write amplification.

\subsection{Hardware Baseline on the Reference Platform}
\label{ssec:hwbaseline}

The current reference run
(\texttt{hardware\_baseline\_gtest\_1775806092.json}) yielded the
normalized factors shown in Table~\ref{tab:hwfactors}.

\begin{table}[t]
  \caption{Normalized Hardware Factors -- Reference Platform
           (run 1775806092)}
  \label{tab:hwfactors}
  \centering
  \begin{tabular}{lrl}
    \toprule
    \textbf{Factor} & \textbf{Value} & \textbf{Tier} \\
    \midrule
    $N_{\text{cpu}}$      & 0.679 & medium \\
    $N_{\text{mem}}$      & 0.619 & medium \\
    $N_{\text{seq}}$      & 0.646 & SSD class \\
    $N_{\text{rand}}$     & 1.000 & high (at ceiling) \\
    $N_{\text{vram}}$     & 0.739 & medium \\
    $N_{\text{h2d}}$      & 0.509 & medium \\
    $N_{\text{d2h}}$      & 0.494 & medium \\
    $N_{\text{dispatch}}$ & 0.402 & low--medium \\
    \bottomrule
  \end{tabular}
\end{table}

Random I/O ($N_{\text{rand}} = 1.0$) is fully saturated, meaning OLTP
workloads are not storage-bottlenecked on this platform.
PCIe dispatch latency ($N_{\text{dispatch}} = 0.402$) represents the
largest gap and constitutes the primary hardware constraint for
GPU-heavy workload classes.

\subsection{Efficiency Results (v0)}
\label{ssec:eff-results}

Table~\ref{tab:eff-v0} presents the first efficiency computation
using the class factors derived from the hardware baseline.
All six classes yield $\text{efficiency} > 1.0$, indicating that the
v0 baseline references are conservative.

\begin{table}[t]
  \caption{Efficiency Results -- v0, Reference Platform}
  \label{tab:eff-v0}
  \centering
  \renewcommand{\arraystretch}{1.1}
  \begin{tabular}{L{1.8cm}R{1.2cm}R{1.3cm}R{1.3cm}R{1.2cm}}
    \toprule
    \textbf{Class} &
    $E_c$ & \textbf{Baseline} & \textbf{Measured} & \textbf{Eff.} \\
    \midrule
    OLTP &
      0.744 & 814.5\,M/s & 796.4\,M/s & 1.314 \\
    OLAP &
      0.647 & 242.6\,M/s & 242.6\,M/s & 1.545 \\
    Storage WAL &
      0.736 & 1.058\,k/s & 1.193\,k/s & 1.532 \\
    Vector GPU &
      0.622 & 351.4\,k/s & 548.7\,k/s & 2.510 \\
    LLM GPU &
      0.620 & 15.9\,k/s  & 15.9\,k/s  & 1.613 \\
    Mixed CPU+GPU &
      0.607 & 717.2\,k/s & 717.2\,k/s & 1.647 \\
    \bottomrule
  \end{tabular}
\end{table}

\subsection{Hardware-Neutral Scores}
\label{ssec:hw-neutral}

Table~\ref{tab:hw-neutral} shows \(\text{score\_hw\_neutral}\) for
the four benchmarks for which both a product target and a measured
value exist on the reference platform.
All scores exceed 1.0, confirming above-target performance relative
to hardware capability.

\begin{table}[t]
  \caption{Hardware-Neutral Scores (run 1775806092)}
  \label{tab:hw-neutral}
  \centering
  \begin{tabular}{L{3.2cm}R{1.3cm}R{1.2cm}R{1.4cm}}
    \toprule
    \textbf{Benchmark} &
    \textbf{Target} & $E_c$ & \textbf{Score} \\
    \midrule
    QueryEngineBench/SimpleEval   & 750\,M/s & 0.744 & 1.427 \\
    VectorIndexBench/InsertPlain  & 280\,k/s & 0.622 & 3.150 \\
    BM\_RawWrite\_WAL\_On/8       & 1.0\,k/s & 0.736 & 1.621 \\
    BM\_EmbeddingCache/100k       & ---      & 0.620 & 1.573 \\
    \bottomrule
  \end{tabular}
\end{table}

\subsection{Hardware Factor Correlation Matrix (v0)}
\label{ssec:corr-matrix}

Table~\ref{tab:corr} presents the qualitative expectation matrix
mapping workload classes to hardware factors.
Entries are High~(H), Medium~(M), Low~(L), or negligible~(0).
This matrix constitutes the prior for the v1 statistical calibration
once sufficient paired runs are available.

\begin{table*}[t]
  \caption{Hardware Factor Expectation Matrix (v0) --
    H=High, M=Medium, L=Low, 0=Negligible}
  \label{tab:corr}
  \centering
  \begin{tabular}{L{2.8cm}
                  C{0.7cm}C{0.7cm}C{0.7cm}C{0.7cm}
                  C{0.7cm}C{0.7cm}C{0.7cm}C{0.7cm}}
    \toprule
    \textbf{Class} &
    $N_{\text{cpu}}$ & $N_{\text{mem}}$ & $N_{\text{seq}}$ &
    $N_{\text{rand}}$ & $N_{\text{vram}}$ &
    $N_{\text{h2d}}$ & $N_{\text{d2h}}$ & $N_{\text{disp.}}$ \\
    \midrule
    OLTP         & H & M & L & H & M & 0 & 0 & 0 \\
    OLAP         & H & H & M & L & L & 0 & 0 & 0 \\
    Storage WAL  & M & M & H & H & 0 & 0 & 0 & 0 \\
    Vector GPU   & M & L & 0 & 0 & H & H & M & L \\
    LLM GPU      & L & M & 0 & 0 & H & H & L & M \\
    Mixed CPU+GPU& M & M & 0 & 0 & M & M & M & M \\
    \bottomrule
  \end{tabular}
\end{table*}

\subsection{Raw Benchmark Details}
\label{ssec:raw-benchmarks}

\subsubsection{Allocator Performance}
\label{sssec:alloc}

\themis{} ships a custom allocator (themis-malloc, inspired by
mimalloc~\cite{Leijen2019}).
Table~\ref{tab:alloc} compares wall-clock throughput against the
system allocator on the v1.8.1-rc2 build (Windows x64,
MSVC~Release, benchmark \texttt{bench\_storage\_performance}).

\begin{table}[t]
  \caption{Allocator Throughput (v1.8.1-rc2, \texttt{bench\_storage\_performance})}
  \label{tab:alloc}
  \centering
  \begin{tabular}{lR{1.3cm}R{1.3cm}R{1.1cm}}
    \toprule
    \textbf{Case} & \textbf{System} & \textbf{Themis} & \textbf{Ratio} \\
    \midrule
    Small alloc ($\leq$256 B)  & 15.5\,M/s  & 128.5\,M/s & 8.3$\times$ \\
    Large alloc ($\geq$64 KB)  & 119\,k/s   & 2.49\,M/s  & 20.9$\times$ \\
    Mixed workload             & ---        & 19.6\,M/s  & --- \\
    RCU read (single thread)   & ---        & 919.8\,M/s & --- \\
    \bottomrule
  \end{tabular}
\end{table}

The 8--21$\times$ allocator speedup is consistent with the mimalloc
free-list sharding approach~\cite{Leijen2019} and is a primary
enabler of the high OLTP throughput reported in
Section~\ref{ssec:longitudinal}.

\subsubsection{WAL Scaling with Thread Count}
\label{sssec:wal-scaling}

Table~\ref{tab:wal-scaling} shows write throughput (ops/s) for three
storage patterns as the thread count scales from 1 to 16.
WAL-ON throughput scales approximately linearly up to 8~threads
(+13\% at 8T vs.\ 1T) but degrades $-25\%$ at 16~threads due to
log-group-commit contention.
WAL-OFF, which bypasses the on-disk journal, scales well up to 8T but
also saturates at 16T.
Mixed 80/20 read-write and secondary-index writes exhibit near-linear
scaling through 16~threads.

\begin{table}[t]
  \caption{Storage Hotspot Scaling (\texttt{bench\_hotspots\_micro}, v1.3.4)}
  \label{tab:wal-scaling}
  \centering
  \begin{tabular}{lR{1.0cm}R{1.0cm}R{1.0cm}R{1.0cm}}
    \toprule
    \textbf{Pattern} & \textbf{1T} & \textbf{4T} & \textbf{8T} & \textbf{16T} \\
    \midrule
    WAL~ON (ops/s)       & 283   & 609   & 1\,193  & 1\,546 \\
    WAL~OFF (k\,ops/s)   & 146   & 370   & 508     & 350 \\
    Mixed 80/20 (ops/s)  & 583   & 1\,289 & 2\,534  & 4\,405 \\
    Sec.\ Index (ops/s)  & 281   & 590   & 1\,056  & 1\,990 \\
    \bottomrule
  \end{tabular}
\end{table}

\subsubsection{Vector Distance Kernel Throughput}
\label{sssec:vec-kernels}

Table~\ref{tab:vec-kernels} reports throughput of the SIMD-accelerated
distance kernels (AVX2, \texttt{bench\_hybrid\_vector\_geo}, v1.3.4).
All kernels operate on 32-bit float vectors.
Throughput decreases roughly as $1/d$ for small $d$ (dominated by
loop overhead) and asymptotes toward a bandwidth-limited regime
around $d=512$.
Geo operations (Haversine, bounding-box containment) operate at
much higher throughput because they use scalar arithmetic on
two-dimensional coordinate pairs rather than full $d$-dimensional
dot products.

\begin{table}[t]
  \caption{Distance Kernel Throughput (AVX2, v1.3.4)}
  \label{tab:vec-kernels}
  \centering
  \begin{tabular}{lR{1.3cm}R{1.5cm}R{1.5cm}}
    \toprule
    \textbf{Kernel} & \textbf{Dim~$d$} & \textbf{Latency} & \textbf{Throughput} \\
    \midrule
    Euclidean    & 64    & 42\,ns  & 23.7\,M/s \\
    Euclidean    & 256   & 209\,ns & 4.8\,M/s \\
    Euclidean    & 1024  & 828\,ns & 1.21\,M/s \\
    Cosine       & 64    & 38\,ns  & 26.4\,M/s \\
    Cosine       & 256   & 205\,ns & 4.9\,M/s \\
    Cosine       & 1024  & 827\,ns & 1.21\,M/s \\
    Haversine    & 2     & 3.6\,$\mu$s/100 pts & 28.0\,M\,pts/s \\
    Bbox contain.& 2     & 64\,ns/100 pts  & 1.56\,G\,pts/s \\
    Vec+Geo filt.& 256   & 1.27\,ms/32K pts & 25.8\,M/s \\
    \bottomrule
  \end{tabular}
\end{table}

\subsubsection{Graph Query Optimizer}
\label{sssec:graph-opt}

Table~\ref{tab:graph-opt} shows plan generation and BFS execution
latency on the v1.8.1-rc2 build
(\texttt{bench\_graph\_query\_optimizer}).
Plan generation is sub-microsecond ($\leq 246$\,ns) and is
cache-resident after the first execution.
BFS execution time grows super-linearly with depth (3\,214\,ns at
depth~2 vs.\ 13\,241\,ns at depth~4), reflecting the exponential
fan-out of the traversal frontier.

\begin{table}[t]
  \caption{Graph Query Optimizer (v1.8.1-rc2)}
  \label{tab:graph-opt}
  \centering
  \begin{tabular}{lR{1.5cm}R{1.5cm}}
    \toprule
    \textbf{Benchmark} & \textbf{Latency} & \textbf{Throughput} \\
    \midrule
    PlanGen -- Shortest Path    & 223\,ns & 4.53\,M/s \\
    PlanGen -- k-Hop (k=2)      & 246\,ns & 4.03\,M/s \\
    PlanGen -- With Cache Hit   & 225\,ns & 4.53\,M/s \\
    BFS -- depth 2 (100 nodes)  & 3\,214\,ns & 321\,k/s \\
    BFS -- depth 3 (100 nodes)  & 6\,038\,ns & 161\,k/s \\
    BFS -- depth 4 (100 nodes)  & 13\,241\,ns & 77\,k/s \\
    \bottomrule
  \end{tabular}
\end{table}

\subsection{Regression Overview v1.3.0 $\to$ v1.3.4}
\label{ssec:regression}

Table~\ref{tab:regression} documents the most significant performance
changes between the v1.3.0 baseline and v1.3.4.
Several regressions marked \emph{critical} are attributable to
intentional infrastructure changes (per-test temporary directories,
single-transaction-per-put RocksDB semantics in test harnesses)
rather than to production code path changes.
This distinction is documented in
\texttt{PERFORMANCE\_COMPARISON\_V1.3.0\_VS\_V1.3.3.md} and must be
accounted for when interpreting absolute values.

\begin{table}[t]
  \caption{Regression Summary v1.3.0 $\to$ v1.3.4}
  \label{tab:regression}
  \centering
  \renewcommand{\arraystretch}{1.1}
  \begin{tabular}{L{2.8cm}R{1.2cm}R{1.2cm}R{0.8cm}L{1.3cm}}
    \toprule
    \textbf{Benchmark} & \textbf{v1.3.0} & \textbf{v1.3.4} &
    \textbf{$\Delta$} & \textbf{Cause} \\
    \midrule
    Vec Insert     & 567\,k/s & 351\,k/s & $-38\%$ & infra change \\
    Sec.\ Index    & 1.78\,M/s & 217\,k/s & $-88\%$ & infra change \\
    Query Engine   & 969\,M/s & 815\,M/s & $-16\%$ & governance \\
    Graph Edges    & 1.47\,M/s & 629\,k/s & $-57\%$ & infra change \\
    TS Insert      & 61.0\,M/s & 49.0\,M/s & $-20\%$ & flush policy \\
    Encrypt 1KB    & 255\,k/s & 191\,k/s & $-25\%$ & EVP overhead \\
    Decrypt 256B   & 60\,k/s  & 41\,k/s  & $-32\%$ & EVP overhead \\
    Geo Content API & 5.6\,k/s & 7.2\,k/s & $+29\%$ & \emph{improvement} \\
    \bottomrule
  \end{tabular}
\end{table}

\subsection{CHIMERA Vendor-Neutral Comparison}
\label{ssec:chimera-compare}

The CHIMERA suite exposes \themis{} benchmarks in a vendor-neutral
form and reports anonymised comparison data obtained from independent
reference runs.
Table~\ref{tab:chimera-qps} and Table~\ref{tab:chimera-vec}
summarise the two primary comparison dimensions: query throughput
and vector search latency.
Statistical methodology follows IEEE~Std~2807-2022~\cite{IEEE2807}:
Welch's $t$-test~\cite{Welch1947} for two-sample comparisons,
Mann-Whitney $U$~\cite{Mann1947} for non-parametric checks,
Cohen's $d$~\cite{Cohen1988} for effect sizes, and
Tukey's IQR rule~\cite{Tukey1977} for outlier removal.

\begin{table}[t]
  \caption{CHIMERA Query Throughput Comparison (95\% CI, Welch's $t$)}
  \label{tab:chimera-qps}
  \centering
  \begin{tabular}{lR{1.0cm}R{1.1cm}R{1.2cm}R{1.0cm}R{1.0cm}}
    \toprule
    \textbf{System} & $N$ & \textbf{Mean} & \textbf{Median} &
    \textbf{CI lo} & \textbf{CI hi} \\
    \midrule
    Alpha (ThemisDB) & 29 & 14\,842 & 14\,813 & 14\,604 & 15\,286 \\
    Beta   & 29 & 12\,678 & 12\,789 & 11\,966 & 13\,090 \\
    Gamma  & 28 &  9\,392 &  9\,431 &  8\,677 &  9\,940 \\
    \bottomrule
  \end{tabular}
  \vspace{0.3ex}
  {\footnotesize All values in queries/s. System names anonymised per IEEE Std~2807 neutrality.}
\end{table}

\begin{table}[t]
  \caption{CHIMERA Vector Search P95 Latency Comparison (ms, 95\% CI)}
  \label{tab:chimera-vec}
  \centering
  \begin{tabular}{lR{0.7cm}R{1.0cm}R{1.0cm}R{0.9cm}R{0.9cm}}
    \toprule
    \textbf{System} & $N$ & \textbf{Mean} & \textbf{P95} &
    \textbf{CI lo} & \textbf{CI hi} \\
    \midrule
    Aurora  & 48 & 8.51 & 10.01 & 7.91 & 8.76 \\
    Nexus   & 49 & 9.43 & 12.40 & 9.02 & 10.02 \\
    Quantum & 50 & 7.59 &  9.27 & 7.25 &  7.93 \\
    Vertex  & 48 & 8.77 & 10.98 & 8.34 &  9.23 \\
    Zenith  & 48 & 9.47 & 12.46 & 8.77 & 10.14 \\
    \bottomrule
  \end{tabular}
  \vspace{0.3ex}
  {\footnotesize Lower is better. $N \geq 28$ per IEEE Std~2807; outliers removed by IQR$\times$1.5.}
\end{table}

System Alpha (ThemisDB) leads query throughput by $+17\%$ over
System~Beta and $+58\%$ over System~Gamma.
In vector search latency the spread is narrower: the fastest system
(Quantum, mean 7.59\,ms) beats the slowest (Zenith, mean 9.47\,ms)
by only $24.7\%$, suggesting that HNSW index performance is largely
converging across the vendor landscape at SIFT-1M scale.

The CHIMERA baseline for \themis{} (v1.5.0-dev,
Table~\ref{tab:chimera-baseline}) also records four core
workload throughput and latency metrics measured within the system
under the same IEEE~Std~2807 protocol.

\begin{table}[t]
  \caption{ThemisDB CHIMERA Self-Baseline (v1.5.0-dev)}
  \label{tab:chimera-baseline}
  \centering
  \begin{tabular}{lR{1.5cm}R{1.3cm}R{1.0cm}R{1.0cm}}
    \toprule
    \textbf{Workload} & \textbf{Throughput} & \textbf{Mean} &
    \textbf{P95} & \textbf{P99} \\
    \midrule
    relational\_sort    & 42\,503\,ops/s & 0.024\,ms & 0.023\,ms & 0.034\,ms \\
    vector\_dot\_product & 75\,835\,ops/s & 0.013\,ms & 0.013\,ms & 0.024\,ms \\
    document\_lookup    & 2.96\,M\,ops/s  & 0.18\,$\mu$s & 0.20\,$\mu$s & 0.25\,$\mu$s \\
    graph\_bfs          & 40\,373\,ops/s  & 0.025\,ms & 0.025\,ms & 0.033\,ms \\
    \bottomrule
  \end{tabular}
\end{table}

% ==============================================================================
\section{Analysis and Discussion}
\label{sec:analysis}

\subsection{Root-Cause Analysis of Performance Gaps}
\label{ssec:rootcause}

\textbf{Query Engine (796.4\,M vs.\ 900\,M\,ops/s).}
The 11.5\% gap between v1.3.4 (814.5\,M) and the 900\,M target
widened to 13.4\% in v1.8.2 due to additional governance-framework
instrumentation in the hot evaluation path.
The primary mitigation is selective hot-path profiling to identify
and defer non-critical checks out of the tight inner loop.

\textbf{Secondary Index (254.9\,k vs.\ 1\,M/s).}
The 74.5\% gap is the largest in the system and is caused by
RocksDB write amplification under the current compaction policy.
Level-based compaction with default $l_0\_file\_num\_compaction\_trigger=4$
generates 4--10$\times$ write amplification for random-key workloads.
Switching to a tiered compaction strategy or pre-sorting inserts by
key range reduces this amplification at the cost of read performance.

\textbf{Storage WAL Sustained Write (1.276\,k/s vs.\ 100\,k/s SLO).}
The production SLO ($\geq 100\,\text{k\,ops/s}$ sustained NVMe write)
is not directly covered by the available proxy benchmark
(\texttt{BM\_RawWrite\_WAL\_On/8} at 1.276\,k/s batch of 8), which
measures fsync-commit overhead rather than sustained NVMe throughput.
A dedicated benchmark targeting 64\,KB sequential writes to NVMe is
required to evaluate this SLO fairly.

\subsection{Benchmark Coverage Gap Analysis}
\label{ssec:coverage}

Across the 33 \themis{} modules, we identify five gap categories
covering in total 29 benchmarks (Table~\ref{tab:gap}):

\begin{table}[t]
  \caption{Benchmark Coverage Gap Categories (33 modules)}
  \label{tab:gap}
  \centering
  \begin{tabular}{L{3.4cm}R{0.9cm}L{3.2cm}}
    \toprule
    \textbf{Category} & \textbf{Count} & \textbf{Root cause} \\
    \midrule
    Feature-gated benchmarks & 7 &
      CMake \texttt{THEMIS\_ENABLE\_*} flags \\
    Disabled stubs & 4 &
      Deactivated \texttt{BENCHMARK(...)} registrations \\
    Missing runtime deps & 3 &
      External library / GPU not present \\
    Proxy-only benchmarks & 12 &
      No 1:1 SLO mapping; indirect metrics only \\
    Missing build artifacts & 3 &
      CMake exclusion / optional submodule \\
    \midrule
    \textbf{Total} & \textbf{29} & \\
    \bottomrule
  \end{tabular}
\end{table}

The dominant category is \emph{proxy-only} (12 benchmarks), where
existing throughput numbers exist but do not directly correspond to
the documented SLO metric.
Notable examples are Cache (proxy via \texttt{BM\_EmbeddingCache})
and Analytics (proxy via \texttt{bench\_olap\_performance}, while the
direct \texttt{bench\_olap\_analytics} is disabled pending API
alignment).

\subsection{Meta-Causes and Remediation Roadmap}
\label{ssec:meta}

Three meta-causes explain the majority of coverage gaps:
\begin{enumerate}
  \item \emph{Proxy vs.\ primary benchmarks.}
        Twelve SLO gaps arise because the available benchmark
        measures a related but not equivalent code path.
        Remediation: add at minimum one benchmark per SLO that
        exercises the exact code path referenced in the SLO.
  \item \emph{CMake feature gating.}
        Seven benchmarks are conditionally compiled behind feature
        flags (e.g., \texttt{THEMIS\_ENABLE\_HNSW\_GPU},
        \texttt{THEMIS\_ENABLE\_LLM}).
        Remediation: add an integration CI job that enables all
        optional features on a GPU runner.
  \item \emph{Disabled benchmark registrations.}
        Four benchmark files exist and compile correctly but have
        their \texttt{BENCHMARK()} invocations commented out
        (typically accompanied by timeout or API-alignment comments).
        Remediation: un-gate or replace with a reduced-scope variant
        that terminates within CI time limits.
\end{enumerate}

\subsection{Known Performance Gaps (D-1..D-7)}
\label{ssec:known-gaps}

Table~\ref{tab:known-gaps} collates the seven explicitly documented
performance gaps extracted from Section~35 of the performance
expectations registry.
All seven require targeted engineering effort before the v2.0 release.

\begin{table}[t]
  \caption{Documented Performance Gaps (Section 35 Registry)}
  \label{tab:known-gaps}
  \centering
  \renewcommand{\arraystretch}{1.1}
  \begin{tabular}{C{0.4cm}L{1.5cm}R{1.3cm}R{1.1cm}R{0.8cm}L{0.9cm}}
    \toprule
    \textbf{ID} & \textbf{Module} & \textbf{Current} &
    \textbf{Target} & \textbf{Gap} & \textbf{Prio.} \\
    \midrule
    D-1 & TS Write      & 200\,k\,pts/s  & 500\,k\,pts/s  & $-60\%$ & High \\
    D-2 & Gorilla Dec.  & 400\,MB/s  & 2\,GB/s   & $-80\%$ & High \\
    D-3 & Vec~vs~FAISS~\cite{Johnson2021}  & 351\,k/s & 600\,k/s & $-41\%$ & Med \\
    D-4 & 2PC~vs~TiDB   & 6.4\,k/s   & 15\,k/s   & $-57\%$ & Med \\
    D-5 & Blob Write    & 741\,ops/s  & 100\,k/s  & $-99\%$ & High \\
    D-6 & Concur. CV    & 20.74       & stable    & ---     & Med \\
    D-7 & Q~vs~ClickHouse & 815\,M/s & 1\,200\,M/s & $-32\%$ & Low \\
    \bottomrule
  \end{tabular}
\end{table}

\textbf{D-1 and D-2 (Time-Series).}
The sustained NVMe write target (D-1) is not yet evaluated by the
available benchmark; the reported 61\,M\,pts/s is for in-memory
append only.
The Gorilla decoder throughput gap (D-2) reflects the lack of an
AVX2/AVX-512 decode kernel; the current implementation is scalar.

\textbf{D-5 (Blob Write).}
The 1\,MB blob write path achieves only 741\,ops/s against a target
of 100\,k/s.
Profiling reveals 99\% of wall time is spent in multi-part
upload coordination within the RocksDB blob extension layer.
A direct \texttt{pwrite} bypass for blobs exceeding a configurable
threshold is the primary proposed mitigation.

\textbf{D-6 (Concurrency Stability).}
Under 10-client concurrent write workload the coefficient of variation
of throughput exceeds 20\%, indicating lock-convoy behaviour in the
RocksDB write-group-commit path.
Increasing the write batch budget or adopting a LMAX-disruptor ring
buffer for the commit queue are candidate mitigations.

\subsection{CI Regression Governance}
\label{ssec:ci-gov}

\themis{} enforces three-tier regression governance via the
\texttt{05-quality\_build\_cross-module-performance-regression-ci.yml}
GitHub Actions workflow:

\begin{itemize}
  \item \emph{Minor} ($\geq 5\%$ degradation): tracking issue created
        automatically; does not block merge.
  \item \emph{Major} ($\geq 10\%$ degradation): PR merge blocked until
        a root-cause comment is attached to the blocking benchmark.
  \item \emph{Critical} ($\geq 20\%$ degradation): immediate
        escalation; PR cannot be merged until the regression is
        resolved or explicitly waived by a maintainer.
\end{itemize}

\noindent
The governance framework operates on \emph{score\_hw\_neutral} rather
than raw throughput, preventing false positives when the benchmark
host's hardware capability differs between runs.
Root-cause analysis is mandatory when
$\text{score\_hw\_neutral} < 0.90$ occurs in two consecutive runs
on the same module (Section~\ref{ssec:calibration}).

% ==============================================================================
\section{Threats to Validity}
\label{sec:threats}

\subsection{Internal Validity}
\label{ssec:int-valid}

\textbf{Infrastructure-induced regressions.}
Several benchmarks listed in Table~\ref{tab:regression} show
large apparent regressions between v1.3.0 and v1.3.4 that are
caused by deliberate changes to the test harness (per-test temporary
directories, single-transaction-per-put semantics for correctness
isolation) rather than by production code changes.
These values must not be interpreted as production performance
regressions.
Mitigation: the \texttt{PERFORMANCE\_COMPARISON} documentation file
explicitly flags all harness-induced changes with root-cause notes.

\textbf{WAL-OFF vs.\ WAL-ON semantics.}
The \texttt{BM\_RawWrite\_WAL\_Off} benchmark bypasses the
write-ahead log entirely, yielding throughput values
($\geq$145\,k\,ops/s) that are not representative of any production
workload, which always requires durability.
All durability-relevant comparisons in this paper use WAL-ON
measurements exclusively.

\textbf{In-memory vs.\ persistent time-series.}
The \texttt{TimeseriesBench/InsertTimepoints} benchmark measures
in-memory append throughput (61\,M\,pts/s at v1.8.2).
Persistent, Gorilla-compressed, multi-tier write throughput is
not yet covered by any automated benchmark (Gap D-1) and may be
substantially lower.

\subsection{External Validity}
\label{ssec:ext-valid}

\textbf{Single platform.}
All measurements are taken on a single Windows x64 workstation with
an Intel i9-10900K.
Throughput numbers may not generalise to Linux ARM64, AMD EPYC,
or cloud-native serverless deployments without recalibration.

\textbf{Benchmark representativeness.}
Google Benchmark micro-benchmarks isolate individual code paths under
unrealistic conditions (e.g., warm caches, single query type, no
concurrent background I/O).
CHIMERA macro-workloads (YCSB, TPC-C/H) provide more realistic
coverage but have not yet been fully executed with production data
volumes ($\geq$10\,TB) on the reference hardware.

\textbf{Hardware model accuracy.}
The v0 efficiency model uses heuristic weights derived from
first-principles analysis.
Until 30 or more paired runs across at least three distinct hardware
classes are available, the statistical significance of the model
cannot be assessed.
All v0 scores exceeding 1.0 should therefore be interpreted as
\emph{conservative baselines}, not as validated performance claims.

\subsection{Construct Validity}
\label{ssec:con-valid}

\textbf{Proxy benchmarks.}
Twelve of the 29 benchmark gaps are proxy-only measurements, meaning
the measured code path does not directly correspond to the
documented SLO metric.
The \texttt{score\_hw\_neutral} values derived from proxy benchmarks
carry higher uncertainty and are labeled \emph{proxy} in all tables.

\textbf{CHIMERA anonymisation.}
The anonymised vendor comparison (Section~\ref{ssec:chimera-compare})
uses system aliases to comply with IEEE Std~2807-2022 neutrality
requirements.
The mapping between aliases and actual systems is not disclosed;
readers should treat comparative scores as indicative rather than
definitive.

% ==============================================================================
\section{Related Work}
\label{sec:related}

Multi-model databases have been surveyed by Lu and Holubov{\'a}~\cite{Lu2018},
who classify systems by their native model and secondary model support.
\themis{} extends this landscape by providing fully native graph,
vector, relational, time-series, and document support with shared
MVCC and a unified query language.

Hardware-normalized evaluation has antecedents in HPC benchmarking.
McCalpin's STREAM benchmark~\cite{McCalpin1995} established the
concept of measuring memory bandwidth as a hardware baseline.
Hennessy and Patterson~\cite{Hennessy2019} discuss the growing
importance of domain-specific hardware.
Our efficiency model applies the same normalization philosophy to
database workloads, extending it to GPU-specific factors
($N_{\text{vram}}$, $N_{\text{h2d}}$, $N_{\text{d2h}}$,
$N_{\text{dispatch}}$) that are absent from CPU-centric HPC models.

YCSB~\cite{Cooper2010} provides workload-portable key-value
benchmarking.
TPC-C and TPC-H~\cite{TPCC,TPCH} define transaction and analytical
workloads, respectively.
ANN-Benchmarks~\cite{AnnBenchmarks} standardize vector search
evaluation.
LDBC-SNB~\cite{LDBCSNB} targets graph analytics.
The CHIMERA suite integrated into \themis{} provides a unified
adapter layer over all five standards.
The statistical methodology used by CHIMERA is grounded in
Welch's~$t$-test~\cite{Welch1947}, Mann-Whitney
$U$~\cite{Mann1947}, Cohen's~$d$~\cite{Cohen1988}, and
Tukey's IQR~\cite{Tukey1977}, following the minimum sample-size and
warmup requirements of IEEE~Std~2807-2022~\cite{IEEE2807}.

Approximate nearest-neighbor search using HNSW~\cite{Malkov2018}
underpins the vector search engine.
FAISS~\cite{Johnson2021} provides GPU-accelerated ANN baselines
against which Gap~D-3 is measured.
The custom allocator design follows the mimalloc free-list
sharding approach~\cite{Leijen2019}, yielding 8--21$\times$ speedup
over the system allocator on allocation-intensive code paths.
RocksDB~\cite{Dong2021} is the storage foundation; its
write-amplification characteristics drive Gap~D-5 (blob write)
and the secondary index gap.
Raft~\cite{Ongaro2014} provides consensus; SAGA~\cite{Garcia-Molina1987}
and 2PC~\cite{Gray1978} handle distributed transactions.
Gorilla delta-of-delta encoding~\cite{Pelkonen2015} is used in the
time-series module; Gap~D-2 documents the planned SIMD acceleration
of the decode path~\cite{Lemire2018}.

% ==============================================================================
\section{Conclusions}
\label{sec:conclusion}

We have presented a longitudinal performance evaluation of \themis{}
spanning versions v1.3.0 through v1.8.2, together with a
hardware-normalized efficiency model that produces platform-independent
benchmark scores, a module-level SLO inventory across 28 modules,
detailed raw benchmark tables, a CHIMERA vendor-neutral comparison,
and an analysis of seven known performance gaps with root causes.

Five major findings emerge:
\begin{enumerate}
  \item Graph edge operations (+87\%) and time-series insert
        throughput (+25\%) both surpassed their product targets in
        v1.8.2, demonstrating effective optimization of data-structure
        hot paths.
  \item PCIe host-to-device transfer ($N_{\text{h2d}} = 0.509$) and
        CPU dispatch latency ($N_{\text{dispatch}} = 0.402$) are the
        dominant hardware bottlenecks for GPU-heavy workloads on the
        reference platform.
  \item All v0 efficiency values exceed 1.0, confirming that the
        initial baseline references are conservative.
        Recalibration to v1 statistical rules requires 30+ paired
        runs across at least three distinct hardware classes.
  \item The CHIMERA query throughput comparison shows \themis{}
        (System Alpha) leading two reference systems by $+17\%$ and
        $+58\%$ respectively; vector search P95 latency is within
        $25\%$ of the field leader.
  \item Seven documented performance gaps (D-1..D-7) remain open,
        of which three (D-1, D-2, D-5) are rated high priority for
        the v2.0 roadmap; all require production-path benchmarks
        that are currently missing.
\end{enumerate}

Future work will focus on (i)~closing the 29~identified benchmark
coverage gaps, particularly the 12~proxy-only benchmarks and 7~high-
priority known gaps; (ii)~running the CHIMERA suite on heterogeneous
hardware (ARM, AMD, NVMe RAID) to build the v1 calibration dataset;
(iii)~extending the efficiency model to encompass network I/O and PCIe
topology for distributed deployments; and (iv)~implementing the
AVX2 Gorilla decode kernel (Gap~D-2) and the LMAX-disruptor write
queue (Gap~D-6).

% ==============================================================================
\balance
\bibliographystyle{IEEEtran}
\bibliography{references}

\end{document}
